
\vspace{1\linewidth}
\addcontentsline{toc}{section}{\cabstractname}
\begin{cabstract}
\textbf{摘要}：随着电力、能源、交通及大型工业设施不断向规模化、复杂化方向发展，设备巡检面临作业环境复杂、人员作业风险高以及设备状态信息获取困难等问题。无人机具有机动性强、可达性高和传感器搭载能力强等特点，已逐渐应用于桥梁、输电线路、光伏场站、风电设施及大型工业设备巡检。研究复杂工业场景下的无人机自主巡检，对于提高设备巡检效率、降低人工巡检风险以及提升异常发现能力具有重要意义。现有无人机自主巡检研究主要涉及环境感知与多模态信息融合、设备缺陷与异常检测、自主巡检任务规划以及多无人机协同等方面。本文分析复杂工业场景下无人机巡检的环境特性和任务需求，系统梳理环境感知、多模态信息融合以及工业设备缺陷与异常检测的研究进展，并进一步总结自主巡检任务规划、动态路径重规划和多无人机协同等相关研究。针对复杂工业场景中传感器退化、异常样本不足、巡检任务动态变化以及真实工业环境验证不足等问题，归纳现有研究的主要局限，并从感知、认知与决策协同的角度提出未来研究思路，为复杂工业场景无人机自主巡检关键技术的研究提供参考。



\textbf{关键词}：无人机；自主巡检；多源感知与信息融合；异常检测；路径规划；多无人机协同
\end{cabstract}
